\section*{1. Abstract}
PipeViz is an interactive cycle-accurate pipeline simulation and visualization platform designed to analyze instruction-level execution behavior in modern processors. The system allows users to input C/C++ programs, compile them into AArch64 assembly, and simulate execution across multiple pipeline architectures including static in-order, scoreboarding, Tomasulo-based scheduling, superscalar, and out-of-order execution models. PipeViz automatically detects and visualizes pipeline hazards such as RAW, WAR, WAW, and structural hazards while presenting a cycle-by-cycle execution timeline through an interactive web interface.

In addition to simulation and visualization, PipeViz integrates a large language model (LLM)-based optimization assistant capable of analyzing source code, generated assembly, and pipeline execution traces to provide context-aware optimization recommendations. The system establishes an iterative workflow where users can observe pipeline inefficiencies, apply compiler or code-level optimizations, and compare the resulting execution behavior. Through multiple case studies involving loop unrolling and cache-aware matrix multiplication, we demonstrate how compiler transformations and LLM-guided optimizations influence hazard frequency, pipeline stalls, and instruction-level parallelism. PipeViz serves both as an educational platform for computer architecture and as an experimental framework for exploring the interaction between compiler optimizations and processor micro architecture. All code is open source and can be found here \cite{github-repo} and The "case\_studies" folder inside the "pipeviz" directory contains the results of the case studies conducted for this report can be found at this link \cite{case-studies}.